% GENERATED FILE. Edit canonical lessons, then run npm run generate:books.
% canonical-source-hash: fnv1a64:f928ce44ea99b348
% canonical-lessons: AR-C439-manzar, AR-C439-burj, AR-C439-nusb, AR-C439-istakshafa, AR-C439-istahaba

\chapter{Landmarks}
\label{ch:ar-landmarks}
\hlchaptermodality{439}

\begin{chapteropening}
\textbf{By the end of this chapter:} \emph{I can say a view, a tower, a monument, to explore and to take along.}

Complete the last lesson of chapter 439: I can say a view, a tower, a monument, to explore and to take along.
\end{chapteropening}

\section[\texorpdfstring{\ar{منظر} (manẓar)}{منظر (manẓar)}]{\ar{منظر} (manẓar) — a view; scenery}

\label{lesson:AR-C439-manzar}



\begin{quote}
Before the new one: say the Arabic for to wander around, then the Arabic for a great traveller.
\end{quote}

\subsection*{You'll want to know: \ar{منظر}}
\textbf{\ar{منظر}} — \emph{manẓar} — \textquotedblleft{}a view; scenery\textquotedblright{}.

\begin{quote}
\ar{كان} \ar{المنظر} \ar{من} \ar{الجبل} \ar{جميلا}.
\end{quote}

\begin{quote}
\emph{kāna al-manẓar min al-jabal jamīlan.}
\end{quote}

\begin{quote}
The view from the mountain was beautiful.
\end{quote}

Plural: \ar{مناظر}, \emph{manāẓir}.

\begin{cousinweb}
\textbf{\ar{منظر}} (\emph{manẓar}) is built on the native Arabic root \emph{n-ẓ-r}, \textquotedblleft{}to look\textquotedblright{} (Lane). The pattern \emph{mafʿal} makes it the thing looked at.
\end{cousinweb}

\begin{grammarlens}[title={What you've built}]
One more word for getting around.
\end{grammarlens}

\subsection*{Your turn}
\begin{itemize}
\raggedright
  \item \emph{Say it:} \emph{manẓar}
  \item \emph{Say it:} \emph{manẓar}, once more
  \item \emph{Say it:} \emph{kāna al-manẓar min al-jabal jamīlan.} — The view from the mountain was beautiful.
  \item \emph{Recall:} say the Arabic for to wander around, then the Arabic for a great traveller, then say \emph{manẓar} again
\end{itemize}

\begin{tcolorbox}[breakable,colback=teal!4,colframe=teal!35!black,title={Before you move on}]
What does \ar{منظر} mean? (\textquotedblleft{}A view; scenery\textquotedblright{}.) Say it once more.
\end{tcolorbox}
\section[\texorpdfstring{\ar{برج} (burj)}{برج (burj)}]{\ar{برج} (burj) — a tower}

\label{lesson:AR-C439-burj}



\begin{quote}
Before the new one: say the Arabic for a great traveller, then the Arabic for a view.
\end{quote}

\subsection*{You'll want to know: \ar{برج}}
\textbf{\ar{برج}} — \emph{burj} — \textquotedblleft{}a tower\textquotedblright{}.

\begin{quote}
\ar{صعدنا} \ar{البرج} \ar{ورأينا} \ar{المدينة} \ar{كلها}.
\end{quote}

\begin{quote}
\emph{ṣaʿidnā al-burj wa-raʾaynā al-madīna kullahā.}
\end{quote}

\begin{quote}
We climbed the tower and saw the whole city.
\end{quote}

Plural: \ar{أبراج}, \emph{abrāj}.

\begin{cousinweb}
\textbf{\ar{برج}} (\emph{burj}) was borrowed from Syriac (EtymDB).
\end{cousinweb}

\begin{grammarlens}[title={What you've built}]
One more word for getting around.
\end{grammarlens}

\subsection*{Your turn}
\begin{itemize}
\raggedright
  \item \emph{Say it:} \emph{burj}
  \item \emph{Say it:} \emph{burj}, once more
  \item \emph{Say it:} \emph{ṣaʿidnā al-burj wa-raʾaynā al-madīna kullahā.} — We climbed the tower and saw the whole city.
  \item \emph{Recall:} say the Arabic for a great traveller, then the Arabic for a view, then say \emph{burj} again
\end{itemize}

\begin{tcolorbox}[breakable,colback=teal!4,colframe=teal!35!black,title={Before you move on}]
What does \ar{برج} mean? (\textquotedblleft{}A tower\textquotedblright{}.) Say it once more.
\end{tcolorbox}
\section[\texorpdfstring{\ar{نصب} (nuṣb)}{نصب (nuṣb)}]{\ar{نصب} (nuṣb) — a monument}

\label{lesson:AR-C439-nusb}



\begin{quote}
Before the new one: say the Arabic for a view, then the Arabic for a tower.
\end{quote}

\subsection*{You'll want to know: \ar{نصب}}
\textbf{\ar{نصب}} — \emph{nuṣb} — \textquotedblleft{}a monument\textquotedblright{}.

\begin{quote}
\ar{زرنا} \ar{النصب} \ar{في} \ar{ساحة} \ar{المدينة}.
\end{quote}

\begin{quote}
\emph{zurnā an-nuṣb fī sāḥat al-madīna.}
\end{quote}

\begin{quote}
We visited the monument in the city square.
\end{quote}

\begin{cousinweb}
\textbf{\ar{نصب}} (\emph{nuṣb}) is built on the native Arabic root \emph{n-ṣ-b}, \textquotedblleft{}to set up, to erect\textquotedblright{} (Lane). A \emph{nuṣb} is a thing set upright.
\end{cousinweb}

\begin{grammarlens}[title={What you've built}]
One more word for getting around.
\end{grammarlens}

\subsection*{Your turn}
\begin{itemize}
\raggedright
  \item \emph{Say it:} \emph{nuṣb}
  \item \emph{Say it:} \emph{nuṣb}, once more
  \item \emph{Say it:} \emph{zurnā an-nuṣb fī sāḥat al-madīna.} — We visited the monument in the city square.
  \item \emph{Recall:} say the Arabic for a view, then the Arabic for a tower, then say \emph{nuṣb} again
\end{itemize}

\begin{tcolorbox}[breakable,colback=teal!4,colframe=teal!35!black,title={Before you move on}]
What does \ar{نصب} mean? (\textquotedblleft{}A monument\textquotedblright{}.) Say it once more.
\end{tcolorbox}
\section[\texorpdfstring{\ar{استكشف} (istakshafa)}{استكشف (istakshafa)}]{\ar{استكشف} (istakshafa) — to explore}

\label{lesson:AR-C439-istakshafa}



\begin{quote}
Before the new one: say the Arabic for a tower, then the Arabic for a monument.
\end{quote}

\subsection*{You'll want to know: \ar{استكشف}}
\textbf{\ar{استكشف}} — \emph{istakshafa} — \textquotedblleft{}to explore\textquotedblright{}.

\begin{quote}
\ar{استكشفنا} \ar{الكهف} \ar{مع} \ar{مرشد}.
\end{quote}

\begin{quote}
\emph{istakshafnā al-kahf maʿa murshid.}
\end{quote}

\begin{quote}
We explored the cave with a guide.
\end{quote}

Present: \ar{يستكشف}, \emph{yastakshif}.

\begin{cousinweb}
\textbf{\ar{استكشف}} (\emph{istakshafa}) is built on the native Arabic root \emph{k-sh-f}, \textquotedblleft{}to uncover\textquotedblright{} (Lane). It is Form X, \textquotedblleft{}to seek to uncover\textquotedblright{}.
\end{cousinweb}

\begin{grammarlens}[title={What you've built}]
One more word for getting around.
\end{grammarlens}

\subsection*{Your turn}
\begin{itemize}
\raggedright
  \item \emph{Say it:} \emph{istakshafa}
  \item \emph{Say it:} \emph{istakshafa}, once more
  \item \emph{Say it:} \emph{istakshafnā al-kahf maʿa murshid.} — We explored the cave with a guide.
  \item \emph{Recall:} say the Arabic for a tower, then the Arabic for a monument, then say \emph{istakshafa} again
\end{itemize}

\begin{tcolorbox}[breakable,colback=teal!4,colframe=teal!35!black,title={Before you move on}]
What does \ar{استكشف} mean? (\textquotedblleft{}To explore\textquotedblright{}.) Say it once more.
\end{tcolorbox}
\section[\texorpdfstring{\ar{اصطحب} (iṣṭaḥaba)}{اصطحب (iṣṭaḥaba)}]{\ar{اصطحب} (iṣṭaḥaba) — to take along; to bring along}

\label{lesson:AR-C439-istahaba}



\begin{quote}
Before the new one: say the Arabic for a monument, then the Arabic for to explore.
\end{quote}

\subsection*{You'll want to know: \ar{اصطحب}}
\textbf{\ar{اصطحب}} — \emph{iṣṭaḥaba} — \textquotedblleft{}to take along; to bring along\textquotedblright{}.

\begin{quote}
\ar{اصطحبنا} \ar{الأطفال} \ar{إلى} \ar{المتحف}.
\end{quote}

\begin{quote}
\emph{iṣṭaḥabnā al-aṭfāl ilā al-matḥaf.}
\end{quote}

\begin{quote}
We took the children along to the museum.
\end{quote}

Present: \ar{يصطحب}, \emph{yaṣṭaḥib}.

Form VIII turns its t into ṭ after ṣ: \emph{iṣṭaḥaba}, not \emph{iṣtaḥaba}.

\begin{cousinweb}
\textbf{\ar{اصطحب}} (\emph{iṣṭaḥaba}) is built on the native Arabic root \emph{ṣ-ḥ-b}, \textquotedblleft{}to be a companion\textquotedblright{} (Lane). It is Form VIII; a \emph{ṣāḥib} is a companion.
\end{cousinweb}

\begin{grammarlens}[title={What you've built}]
One more word for getting around.
\end{grammarlens}

\subsection*{Your turn}
\begin{itemize}
\raggedright
  \item \emph{Say it:} \emph{iṣṭaḥaba}
  \item \emph{Say it:} \emph{iṣṭaḥaba}, once more
  \item \emph{Say it:} \emph{iṣṭaḥabnā al-aṭfāl ilā al-matḥaf.} — We took the children along to the museum.
  \item \emph{Recall:} say the Arabic for a monument, then the Arabic for to explore, then say \emph{iṣṭaḥaba} again
\end{itemize}

\begin{tcolorbox}[breakable,colback=teal!4,colframe=teal!35!black,title={Before you move on}]
What does \ar{اصطحب} mean? (\textquotedblleft{}To take along; to bring along\textquotedblright{}.) Say it once more.
\end{tcolorbox}
